\silentchapter{2 Ikkje alle posisjonar er like sannsynlege}{ikkje alle posisjonar er like sannsynlege}
\prop{2}{}{Ikkje alle posisjonar er like sannsynlege\footnote{Sannsynet er vekta av den samla systemtilstanden; ingen tenestekonfigurasjon flyt fritt.}.}
\prop{2.1}{d}{Eit \textbf{seleksjonstrykk\footnote{Eit trykk kan vere prosessorkostnad, brukaradopsjon, eller overføringsbandbreidde.}} er ein faktor som endrar sannsynsfordelinga over formrommet.\footnote{Wright (1932)}}
\prop{2.11}{t}{Det er ikkje ein regel som dikterer éi form, men ein gradient som gjer visse posisjonar meir sannsynlege enn andre\footnote{Gradienten tvingar fram eit mønster utan å diktere den eksakte kodelinja eller grensesnitt-pikselen.}.}
\prop{2.111}{t}{Eit seleksjonstrykk åleine avgjer ikkje ein posisjon.}
\prop{2.1111}{t}{Det avgrensar kvar posisjonen ikkje kan vere.}
\prop{2.11111}{t}{Form er det som vert att når alle seleksjonstrykk har utelukka det dei kan utelukke\footnote{Form er restverdien etter at alle optimaliseringskrav har barbert vekk dei ikkje-levedyktige alternativa.}.}
\prop{2.11112}{t}{Ho er ein rest, ikkje ein konstruksjon\footnote{Systemdesign er soleis ein reduktiv prosess av feilretting, ikkje ein additiv prosess av skaping.}.}
\prop{2.111121}{t}{Eit trykk som ikkje utelukkar nokon region er ingen trykk; det er ein tom kategori.}
\prop{2.111122}{t}{Eit trykk som utelukkar alle regionar unntatt éin er ingen trykk; det er ein determinisme.}
\prop{2.111123}{t}{Dei trykka vi observerer ligg mellom desse ytterpunkta.}
\prop{2.12}{a}{For kvar funksjonell klasse verkar det alltid meir enn eitt seleksjonstrykk samstundes\footnote{Eit system i stabil drift er alltid fanga mellom kryssande trykk, som kravet til tryggleik versus kravet til yting.}.}
\prop{2.121}{t}{Einsidige forklaringar av form er alltid ufullstendige\footnote{Å tilskrive suksessen til eit tenestesystem utelukkande til brukarvenlegheit er ein logisk forenkling som overser teknisk og økonomisk trykk.}.}
\prop{2.1211}{t}{Dersom berre eitt trykk verka, ville alle former i klassen vere identiske eller tilfeldig fordelte langs éin akse.}
\prop{2.12111}{t}{Den observerte ikkje-uniformiteten med variasjon langs fleire uavhengige aksar krev fleire uavhengige trykk.}
\prop{2.1212}{t}{At eit trykk er latent tyder ikkje at det er fråverande.}
\prop{2.12121}{o}{Det er vanskeleg å telje trykk som ikkje varierer i det observerte utvalet.}
\prop{2.121211}{t}{Metodisk krev dette variasjon i vilkåra, ikkje berre variasjon i formene.}
\prop{2.13}{a}{For minst eitt par av seleksjonstrykk peikar gradientane i motstridande retningar i minst éin region av formrommet\footnote{Spenninga mellom dataminimering og rik funksjonalitet peikar i eksakt motsette retningar av formrommet.}.}
\prop{2.131}{t}{Av 2.12 og 2.13: sidan fleire uavhengige trykk verkar og dreg i ulike retningar, kan ikkje alle verte optimalt tilfredsstilte samstundes.}
\prop{2.1311}{t}{Kvar realisert form er eit kompromiss.\footnote{Michl (1995)\footnote{Eit programvarekompromiss manifestet som ein mikroteknestearkitektur; ho tolererer bandbreiddetap for å vinne agens-isolasjon.}}}
\prop{2.13111}{t}{Eit kompromiss er ikkje ein svakheit; det er den einaste moglege balansen under dei vilkåra som rådde.}
\prop{2.13112}{t}{Sidan det finst fleire gyldige kompromiss, er skilnaden mellom to former ikkje ein feil, men eit uttrykk for at kreftene kan balanserast på meir enn éin måte\footnote{Ulike operativsystem representerer ulike gyldige kompromiss for det same overordna settet av seleksjonstrykk.}.}
\prop{2.13113}{t}{Å kalle ein form «suboptimal» føresett at ein veit kva ho er suboptimal i høve til.}
\prop{2.131131}{t}{Utan eksplisitt spesifisering av trykk-settet er «suboptimal» ikkje falsifiserbart\footnote{Uten å matematisk vekte yting mot lesbarheit er vurderinga av eit kode-paradigme uformell.}.}
\prop{2.14}{t}{Funksjonen definerer klassen og med det formrommet.}
\prop{2.141}{t}{Innanfor klassen er funksjonen konstant\footnote{Den underliggande forretningslogikken gjev ingen informasjon om kvifor to klientar med ulikt grensesnitt løyser oppgåva ulikt.}.}
\prop{2.1411}{t}{Det som er konstant har null varians og ber null informasjon om kvifor éi bestemt form vart realisert framfor ei anna.}
\prop{2.14111}{t}{All observert formvariasjon under konstant funksjon er driven av dei andre seleksjonstrykka.}
\prop{2.141111}{i}{Stolen og sofaen har ulik funksjon og utgjer ulike klasser.}
\prop{2.141112}{i}{To stolar med ulik stil har same funksjon. All skilnaden mellom dei er ikkje-funksjonell.}
\prop{2.141113}{o}{At materialet er den faktiske mekaniske avgrensninga, medan stilperioden ber ingen ergonomisk bodskap, og likevel forklarer stilperioden meir formvariasjon enn materialet, er ein direkte empirisk motbevising av den funksjonalistiske doktrinen innanfor denne klassen.}
\prop{2.2}{d}{\textbf{Materialaffordansen\footnote{Affordansen i eit grafisk grensesnitt er ikkje ein farge, men kva for rekke av hendingar fargen logisk opnar for å initiere.}} er dei operasjonane eit materiale tilbyr og motset seg under gjevne kostnads- og tilgjengelegheitsavgrensingar.\footnote{Gibson (1979)}}
\prop{2.21}{t}{Signaturen er probabilistisk: repertoaret bind sannsynsfordelinga, ikkje det einskilde objektet\footnote{Ein API-dokumentasjon er ei skildring av dette latente repertoaret av moglege tilstandar.}.}
\prop{2.211}{o}{Signaturen er latent.}
\prop{2.2111}{o}{Det teoretiske repertoaret slår ikkje nødvendigvis ut i observerbar geometri.}
\prop{2.21111}{o}{Materialet er til stades, men algebraen manglar\footnote{Kapasiteten for distribuert konsensus låg i kryptografi og nettverksprotokollar i tiår før blokkjedene tok posisjonen.}.}
\prop{2.21112}{o}{Det som endrar seg over tid er ikkje materialet, men kva operasjonar som er praktisk tilgjengelege.}
\prop{2.211121}{i}{Latensen kan vere lang. Stål var tilgjengeleg i hundreår før den fyrste røyrstolen; materialet fanst, men operasjonen som kunne utnytte det, fanst ikkje.}
\prop{2.211122}{t}{Formhistoria er difor systematisk langsamare enn materialhistoria\footnote{Systemarkitekturen vil alltid henge etter maskinvarens affordansar, av di operasjonane tek tid å formalisere til grammatikk.}.}
\prop{2.211123}{o}{Nye materiale arvar formspråket til det substratet dei erstattar, heilt til eigne affordansar pressar dei ut i nye regionar av formrommet.}
\prop{2.22}{t}{Ein samlevariabel korrelert med alle samtidige seleksjonstrykk bør predikere formvariasjon betre enn kvart einskildtrykk.}
\prop{2.221}{o}{Stilperiode er ein slik samlevariabel.}
\prop{2.2211}{t}{Han er ikkje ein forklaring; han er ein proxy som absorberer alle samtidige trykk i éin etikett\footnote{Termar som 'Web 2.0' eller 'skya' skildrar kvar vandringa stansa, men forklarar ingen av dei underliggande kausale drivkreftene.}.}
\prop{2.22111}{t}{Stilkategoriane skildrar kvar vandringa stansa. Dei forklarer ikkje kvifor ho stansa der.}
\prop{2.22112}{t}{Å blande saman den deskriptive og den kausale krafta til stilomgrepet er den vanlegaste feilen i formhistorieskrivinga.}
\prop{2.221121}{t}{Stilperioden er eit symptom; seleksjonstrykka er årsaka.}

